\chapter[Sermon 1. Of the Author of the Book of Apocalypse, of the Argument and Parts Thereof: Finally of the Sundry Use, and Most Profitable Commodity of the Same.]{Of the Author of the Book of Apocalypse, of the Argument and Parts Thereof: Finally of the Sundry Use, and Most Profitable Commodity of the Same.}
\label{sermon:001}
\markright{Sermon 1. Of the Author of the Book of Apocalypse, of the ...}

\index{Antichrist}\index{Repentance}The prophets of God of the old Testament were God’s messengers to the people, expositors of Moses or God’s law, and even ecclesiastical preachers, who applied the doctrine taken from the law of God to the people and times in which they lived, for the edification of congregations. And they all with one accord have chiefly handled two things in their sermons. For first, they have reproved the corrupt manners of all states in their time, alleging the rule and prescription of God’s law; and exhorting all men to repentance unto God most acceptable. And to many that were uncurable they threatened all kinds of plagues, which they setting forth with all beauty of speech, showed them plainly to be seen with the eye, if they apply them so might be made afraid, and healed. Secondly, they promised and set forth by the mouth of God, the Lord Christ, the true Messiah: Whom also they described lively, and with all his holy church, teaching the faith in Christ, and showing what good things are prepared for the faithful in Christ; and also what be the true duties of piety and godliness. Neither have they concealed such things as were needful to be known of Antichrist; admonishing us most diligently that we should beware of that wolf, or rather the deepest pit of all abominations, and that we should stand fast in the sincere faith of Christ. \&c.

\index{Apocalypse (Book of)}I have undertaken, through God’s help and your prayers, to explain the Apocalypse, which is a doctrine concerning the matters of Christ’s church revealed from heaven in Christ’s glory, a summary of all godly religion, an exposition and brief declaration of the Prophets, and a prophecy of the New Testament and the history of the church.

However, since this book is despised by both good and well-learned men, and because most people are firmly convinced that it is an unprofitable study, I will speak somewhat about this matter here.

\index{Daniel (Book of)}First, many abhor this book for this cause only that it is full of visions, types, and figures, for they suppose this does not become the Evangelical and Apostolic doctrine of the New Testament. But by the same reasoning, a portion of Daniel should be cast away. Nevertheless, this is commended to us by Christ himself in Matthew. The greatest matter of all, namely the calling of the Gentiles, was shown to Saint Peter by a vision, as appears in the Acts.

\index{Jerome, Saint}And the prophet Joel also said how the people of the New Testament should see visions. And so does the blessed Apostle Jerome, expound the same place in the Apocalypse of the Apostles, speaking of the people of the New Testament. And our Savior Christ in the Gospel, proposed and declared to the people many mysteries by parables, and in a manner by feigned fables, as they call them. And how much do you think these visions, types, and figures of John differ from the same? This kind of speaking does not darken matters, but makes them plain; and it contributes much to the efficacy and clarity, and to the firming of the memory. For by this means, matters are not only declared with words and heard with the ears, but they are set forth to be seen by the eyes, and after a sort are fixed in the memory.

Many for this cause attribute much to painting. But I suppose that I may much more rightly attribute very much to this manner of speaking and teaching, where the matter is uttered, not by a colored, dull, and dead painting, but as it were with a speaking and lively manner set forth to be seen by the eyes. Which is therefore proposed, that men might rightly and exactly understand the same. Although therefore this whole book in a manner consists of visions and figures, yet shall we, through the inspiration of God’s grace, show in our exposition that it pertains to the perspicuity and plainness, and not to the obscuring or darkening of most high and godly matters. I will bring my exposition out of the very scriptures, by comparing and joining thereunto the rule of faith and charity. I will search out the circumstances, the things that follow and go before; I will bring similitudes and dissimilitudes; I will add also thereunto the experience of things, and the faith of histories. Which manner of expounding the scriptures all interpreters have always granted to be sound and true.

If better things shall be revealed to others, I will gladly, after the precept of the apostle, give place unto my betters. For I offer these my doings to be weighed by the godly, upon this condition, that they should try all things, and that which they shall find to be good, to hold fast.

\index{John, Saint}\index{Erasmus of Rotterdam}\index{Luther, Martin}\index{Eusebius of Caesarea}\index{Arethas of Caesarea}\index{Papias}Secondly, they object that both new and old men of no small authority have doubted this book and its author, and have also scorned it as full of fables and unworthy to be reckoned canonical. Let those who think so give me the same liberty I desire of them, which they themselves usurp and deem lawful. For if the book of the Apocalypse should therefore seem worthy of scorn because some notable men, both old and new, have doubted of its authority, why may it not regain its authority if I show that the best doctors of the church, both old and new, have held a right good opinion of this book? And here, to the intent that I would dissemble nothing at all, I am not ignorant that Doctor Martin Luther, a man notably learned, in his first edition of the New Testament in Dutch, with a sharp preface set before, has rejected this book as it were with a dagger. However, it offended good and well-learned men that he held this judgment, which they found lacked both wit and modesty. The same man, therefore, weighing all things more uprightly and diligently, when he corrected his Dutch Bible… My worshipful master also seems not to have set very much by his book, and to have ascribed it not to John the Apostle, but to John who they called a divine. But herein there is no doubt but that he followed plainly Erasmus of Rotterdam, who in his annotations upon the New Testament, in all the Greek copies, which I have seen, states that the title was not of John the Apostle, but of John the divine. Erasmus adds that amongst the Greeks and certain old writers, men doubted of this author, which thing he declares by the testimonies of Eusebius and Jerome, whose opinion shall be spoken of straightways. But the exemplar or Spanish copy which is set forth after the faith of the most ancient and approved Greeks, exhibits to us such a title of this book: That is, the Apocalypse of the holy Apostle and Evangelist St. John Divine. For the ancient writers say how St. John the Apostle and Evangelist, for his excellent writing of the Son of God, was commonly called Divine. Whereof it follows that this title does attribute to, and not take from, St. John this book. Certainly, Arethas was also a Greek author, a bishop of Caesarea: Of the ancients, he says, certain have plucked this Apocalypse from the tongue of that well-beloved John, ascribing it to another; but it is not so. For that same Gregory, who as well as he is called a Divine, accounts this amongst those scriptures which utterly lack a suspicion of counterfeiting, saying, as the Apocalypse of St. John teaches me. And the same man a little after: But if this book was written by the mouth of the Holy Ghost, St. Basil, Cyril, Papias and Hippolytus, fathers of the church, are men to be credited. Thus saith he. What shall we say that Erasmus confesses that the consent of the world, and the authority of the church, are of such force with him that he dare not refuse this book?

\index{Oecolampadius, Johannes}Let us hear now the judgment of that most excellent man, Dr. John Oecolampadius, the most faithful pillar of the church of Basel, and excellently learned in prophetical matters and in all the canonical scriptures, concerning this book, which he left us written in the twelfth chapter of the second book of his commentaries upon Daniel. But St. Jerome, the paraphraser or expounder of the Prophets, says he—[?how much he attributes to this our author]—whom I marvel why certain with so rash a judgment do reject as a dreamer, and frantic, and an unprofitable writer to the church. Nevertheless, he propounds and sets forth very many of the most secret and hidden things of the old testament and of the Prophets. But those great men do betray what a sense of their own importance they have. Whose judgments I would verily rather scorn as profane than I would cast away such a treasure. I could here bring forth goodly testimonies of other new writers, but that I make haste to the judgments of the ancient fathers.

\index{Justin Martyr}\index{Irenaeus}\index{Tertullian}The earliest of all writers after the Apostles, Justin and Irenaeus, the noble Martyrs of Christ, attribute this book to John the Apostle. For Eusebius, in chapter eighteen of the fourth book of the \textit{Ecclesiastical History}, states that Justin mentions the Apocalypse of John, plainly saying that it is the Apostle’s. Jerome also writes in the life of blessed Justin that Justin explained the Apocalypse of Saint John, but that explanation does not survive to any great extent as far as I know. The same author writes that Irenaeus presented the Apocalypse of Saint John with a commentary, which is also unavailable. He himself, said to have lived around the year of our Lord 160, testifies plainly in the fifth book against the Valentinians that this revelation was shown to John the Apostle a little before his time. We cite certain words of his in chapter thirteen of this book. Tertullian, who lived around the year of our Lord two hundred and twenty, in the fourth book against Marcion, although he says that Marcion rejects the Apocalypse of John, nevertheless states that the order of bishops reckoned up to the very beginning affirms Saint John to be its author. In serious matters and when reasoning against heretics, he readily uses the testimonies of this book.

\index{Cyprian, Saint}\index{Origen}The same things are also recorded concerning the blessed Martyr Saint Cyprian, under the name of John the Apostle, in his Epistles, treatises, and Sermons. Eusebius also, in chapter 18 of the fifth book of the \textit{Ecclesiastical History}, shows that Apollonius, a very ancient writer, uses the testimonies of the Apocalypse of Saint John. And likewise Theophilus, Bishop of Antioch, who affirms this in chapter 23 of the fourth book of the \textit{Ecclesiastical History}. Also Origen, a great man in the church of God, in chapter 25 of the sixth book of the same Eusebius. He wrote, says he, the Apocalypse, which rested upon the Lord’s breast, and so forth.

I have thus far recounted the opinions of the most ancient Martyrs and Doctors of the Christian church, concerning the Apocalypse, meaning Justin, Irenaeus, Tertullian, Cyprian, Apollonius, Theophilus, and Origen. I will shortly bring yet more judgments, both of Greek and Latin writers, of most authority in the church, agreeing with the minds of those that we have alleged hitherto. However, I will first briefly touch upon such things as Dionysius of Alexandria left written about the same book in the twenty-fifth chapter of the seventh book of Eusebius, whom I suppose they have followed, as many after him have spoken against this book. He says that diverse of his predecessors did utterly reject and reject this book. Neither does he hide the cause why they so did, for that the kingdom of Christ is affirmed therein to be earthly. Whereunto doubtless they referred that precious city, and the rest which under terrestrial kinds, figured spiritual things. Which when we, in the treating thereof, have dissolved, declaring that it does not edify the earthly kingdom of Christ, but a spiritual and celestial one, no man, I think, will reject a good and godly book, because certain abusing the testimonies thereof, give unto it a wrong sense.

\index{Millenarianism}Have heretics twisted many passages of scripture to defend their errors, should the authority of scripture itself then be called into question? I favor neither the Chiliasts nor the Millenaries in regard to this book; I give them no support.

\index{Resurrection}Eusebius speaks very well at the end of the third book, when speaking of Papias, the first author of the Millenaries. He thought, says he, that after the resurrection Christ should reign here corporally with his thousand years on earth. Which I suppose he thought because he did not understand well the Apostles' words, nor did he consider well those things that were spoken of him under figures, for that he was endowed with a small judgment [? understanding].

But in the meantime, Dionysius himself, I say, would not reject this book. He adds immediately that he thinks it not yet to be the book of John the Apostle, but of someone else, but who that someone should be, he did not know. He gathers also by certain conjectures, by the phrase of speech, and handling of the book, and by the unlikeness of wit, that this book should be another’s than he who wrote the Gospel and Epistle. But seeing that the arguments of the story and Epistle are so diverse, that neither are like, and the argument of the book of Revelation most diverse of all: why should it seem marvelous, though it agrees not with them in all things?

\index{Hebrews, Epistle to the}No one can deny that there is great agreement in doctrine. The Epistle to the Hebrews appeared to many to favor the Novatians or Catharites in chapter six and ten. The difference in style was noted to differ from the rest of Saint Paul’s Epistles. But if we were to judge the holy scriptures in this way, I do not know what would be firm and sure enough. Leaving this disputation in suspense, I will now proceed to bring forth the judgments of other old writers concerning this book.

\index{Patmos}Eusebius, nicknamed Pamphilus, Bishop of Caesarea, living during the reign of Emperor Constantine, and a devoted student of ancient writers, whom many suspect sought to undermine the authority of this book to their own advantage, eloquently affirms in chapter eighty of the third book of his history that John was exiled to Patmos and wrote his Apocalypse there, while denouncing the tyranny of Domitian.

And other historians do likewise. He again, in the twenty-third chapter of the third book, concerning the Apocalypse, says that the opinions of people are diverse, some approving and others reproving it. Again, when he should present his opinion concerning the canon of the New Testament in the twenty-fifth chapter, he joins the Apocalypse with the books of unquestioned authority, although he does not dissemble that he will show elsewhere what other people think of it. While he performs this, he incorporates many more and better, who judged the Apocalypse to be of Saint John the Apostle, and embraced it as a most godly book, than those who denied or reproved it.

\index{Epiphanius of Salamis}Epiphanius, Bishop of Salamis in Cyprus, a Greek author, also clearly ascribes this book to John the Apostle. Read what he has written against heretics in the fifteenth heresy. And Jerome attributes much to this Epiphanius.

\index{Augustine, Saint}\index{Ambrose, Saint}Jerome himself ascribes this book to John the Apostle, to Paulinus: “The Apocalypse of John,” he says, “has so many sacraments as it has words.” Moreover, Philastrius, Bishop of Griria, whom Augustine says he saw with Ambrose at Milan, considers them heretics who reject the Apocalypse of John, and say that it is not of John the Apostle, but of Cerinthus, an heretic. Ambrose himself alleges in his books testimonies from the Apocalypse, under the name of John the Apostle.

\index{Primasius}Augustine embraced this book as Apostolic and referred the same to his church, leaving certain treatises upon the same. Primasius, also Bishop of Utica in Africa, explained it as Apostolic. Of Bed and the rest of that sort, I speak nothing, since his opinion is known to all men. Andreas, also Bishop of Caesarea, wrote upon this book, as Arethas reports in his commentaries, an opinion I declared before.

I believe I have sufficiently confirmed the authority of this book, against those who diminish it. But the strongest evidence seems to be the book itself and the way it is handled, proving that it proceeded from the Apostle. This we shall demonstrate in the treatise itself. But if those blessed fathers in their time did interpret the Apocalypse for their churches, why should it not be lawful for us also in our time to interpret it for our people, who live at the end of the world, where all things are now more fully realized than they were in their time? Indeed, these things serve most chiefly for us and for our time, especially as we travel and are exercised under Antichrist.

Therefore, it is in vain for many to claim that this book is obscure and cannot be understood, and for that very reason should not be read in the church without any profit or benefit.

For to speak nothing of this, nothing is found in holy scripture which does not have an excellent fruit. Neither must we immediately despair of true understanding, even though at first sight the holy scripture seems obscure, which is opened by God himself, and is obtained through prayer and godly exercises. Certainly, we are not ignorant that many would prefer that nothing be spoken of Antichrist, so that he might reign more carelessly, and they themselves be less subject to peril. But Christ commands us to trouble him. Let us therefore proceed in the work of the Lord.

And where it offends them, that John makes little mention or none of Christ, whereas the manner of the Apostles is always to intimate Christ and the grace of redemption: we suppose this same book will be looked upon more thoroughly to prove the contrary. Whose argument now I will recite.

\index{Zechariah (Book of)}The prophet Zechariah in chapter three presents the entire mystery of Christ to all people’s eyes in a most evident figure. He sees Jesus, the high priest, clothed in shameful garments, and like a coal taken from the fire, suffering much opposition from the Devil. Immediately afterward, he sees the same person cast off the shameful clothing and put on white garments, to be glorified, and proclaimed king, priest, and savior of all.

This is what the Apostle and Evangelist John explains. First, after presenting the Gospel, he describes Christ in humble clothing, detailing the great opposition he endured from the wicked until he was nailed to the cross. He also touches upon his glory there. However, the Apocalypse appended to that Gospel expands upon this, showing him in a white garment and in glory, truly demonstrating how, after this humbling, he was exalted and obtained a name above all names. And now, being in glory, he nevertheless works in the church, the savior of all the faithful within the church. In his Epistle, he praises this entire mystery of piety and urges it upon all people.

The entire book is divided into six sections.

For first is set the title with the beginning and some of the work, and with a brief narration: and all this in the first part of the first chapter.

Secondly, from the midst of the first chapter to the fourth chapter, Christ’s reigning in glory is described, at the right hand of the Father, and it is declared how he is active in the Church through his Spirit and the ministry of his word. What things he teaches from heaven, and what the sincere doctrine of the church is, are revealed. Also, the restoration of fallen churches and their preservation are shown.

3 From chapter four through twelve, Christ continues to earnestly warn his Church through seven seals and seven trumpets, describing the events that will befall the Church, all of which are most justly ordained by God himself through the Lamb, Christ.

4 Moreover, from chapter twelve to fifteen, the conflict of the Church with the old Serpent and with the old and new beast is more fully described. Here, too, that monstrous Tyranny, both old and new, and very Antichrist himself is clearly portrayed in his colors; nevertheless, these things are afterward also again more plainly declared.

5 And from the fifteenth Chapter to the twenty-second Chapter, are recited the pains and torments of Antichrist and Antichristians, and the destruction of the same and the condemnation of all the wicked. Also the Judge, Christ, is set forth, and the process of an external judgment is figured. There is also rehearsed the triumph, joy, and reward of Saints. Where also heaven itself is opened to be seen by our eyes, that now we may by faith look into the same. The depth of Hell is opened, that we may look into it also: and may take good heed that we be not thrown thither headlong.

Finally, at the end of chapter twenty-two, follows the conclusion and commendation of the work, with its sealing.

And here I will not conceal another division of this work, not to be disregarded. I observe that commentators have followed this structure. First, they present the title and beginning. Then, they append the entire work itself, divided into seven visions. Indeed, the number seven is most frequent and, as it were, distinctive to this book. Finally, they add a conclusion that summarizes the work, contained within the last chapter. And these visions are contained within their boundaries.

In the first three chapters, the initial vision is explained, exhibiting Christ to us reigning in glory, governing, ordering, correcting, and preserving his church.

The second vision begins in the fourth chapter and extends to the eighth chapter. This section portrays God himself and his Christ, presenting them for observation, and it commends their most just governance of all things in the world, opening seven seals.

The third vision features seven angels sounding with seven trumpets. This section extends to chapter twelve.

The fourth vision depicts the struggle of the woman with the serpent, and presents to us the old seven-headed beast and the new two-horned beast, which are descriptions of Antichrist, as found in chapters twelve, thirteen, and fourteen.

In the fifth vision, seven angels pour out seven vials of God’s wrath, continuing to the seventeenth chapter.

From there begins the sixth vision, and it extends to chapter twenty-one, discussing God’s just judgment against Babylon, the great prostitute, and against those who oppose Christ, ultimately against all unrepentant and wicked people.

The seventh and final vision presents to the eyes of all the faithful the glory and eternal blessedness of the saints. Indeed, this section of the work possesses great beauty and affinity with the rest of the things which in this book are all, in a manner, treated by the seventh number. Let the reader follow whichever he prefers.

Now concerning these things, every person may perceive that this book is entirely apostolic and exceedingly profitable to us all, especially those whom the end of the world has overtaken. And this book will be easier to understand, for now all things are in a manner accomplished.

Daniel was thought to have told of strange visions when before the monarchies, he prophesied the monarchies. But after those things were accomplished which he prophesied, he seemed unto many to have compiled a fictitious story. The self same, I am sure, you will judge also: this same book of Saint John. A few prophecies only we shall recite.

First we have in this book a most full description of Christ reigning in glory, our king, I say, and Bishop; and how he governs the Church, and is the Savior of all the faithful. We have also a most gallant description of Christ’s Church, and how the same may be built, repaired, and maintained. Then we have a perfect description of Antichrist, of his members, and synagogue of his counsels, crafty devises, kingdom, cruelty, and destructions of the same; from which it bids us beware. Moreover, we have an abridgment of history from Christ’s time unto the world’s end.

Finally, an absolute and certain prophecy of things to come, so that we need not the prophecies of Methodius, Cyril, Merlin, Bridget, Nolhard, and certain triflers.

Furthermore, we have great consolation and comfort for the church in adversity, while we see the Lamb open the Seals, and that all things are accomplished by God’s providence, and that there is an end to evil. And the church shall endure forever, in spite of all the devils in hell. Lastly, we have a most plentiful and sure doctrine of the Judge and final judgment, of punishments and of rewards.

All these things I say, the treatise itself will plainly demonstrate, for our edification through Jesus Christ our Lord.
